\begin{topic}[id=iot_cyberattacks,
             difficulty={Mittel},
             type={Sicherheitsanalyse + Fallbezug},
             prerequisites={Grundlagen IT-Sicherheit; Basiswissen IoT-Systeme hilfreich},
             practice={optional},
             owner={Dozententeam},
             updated={2026-04-09},
             tags={ IoT-Threats, OWASP IoT, Hardening, Updates, Secrets, Logging }]{Cyberangriffe auf IoT-Anwendungen}

\TopicDescription{Das IoT vergrößert die Angriffsfläche: schwache Voreinstellungen, unsichere Update-Wege, mangelhafte Schlüsselverwaltung und falsch konfigurierte Broker/ACLs. Mit den OWASP-IoT-Top-10 strukturierst du typische Risiken und leitest priorisierte Maßnahmen ab: sichere Werkseinstellungen, robuste Update-Mechanismen, Härtung von Standarddiensten und Telemetrie/Logs, die im Feld nutzbar sind. Ziel ist ein kompakter, umsetzbarer Sicherheits-Leitfaden für reale Projekte.}

\TopicLeitfragen{\item Welche OWASP-Risiken treten in der Praxis am häufigsten auf?
\item Welche Maßnahmen liefern früh den größten Sicherheitsgewinn?
\item Wie sieht ein schlanker Logging-/Monitoring-Pfad aus?
}
\TopicScope{\item Keine Pen-Test-Anleitung.
\item Keine proprietären Protokolldetails.
}
\TopicPractice{Demo-Stack mit zwei Fehlkonfigurationen und anschließender Härtung.}

\TopicKeywords{IoT-Threats, OWASP IoT, Hardening, Updates, Secrets, Logging}

\nocite{enisaThreatLandscape,owaspIoTTop10}
\end{topic}
